\begin{lemma}\label{lem:decic242-quotient5}
Put $X=X_1(2,42)$, and let $G_9$ be generated by
\[
 a(P,Q)=(21Q,P+Q),\qquad b(P,Q)=(P,25Q).
\]
Then $G_9\simeq C_3\times C_3$ is defined over $\mathbb Q$.
The quotient $Y=X/G_9$ has genus $5$ and the following function
field. Put
\[
 t=u^2,\quad D=t^2+13t+49,\quad N=t^2+5t+1,\quad
 M=t^4+14t^3+63t^2+70t-7.
\]
Adjoin the coordinates $(x_3,y_3)$ of a nonzero point of order $3$
on
\[
 E_u:\quad y^2=x^3-\frac{N}{48D}x+\frac{M}{864D^2}.
\]
Thus $x_3$ satisfies the $3$-division polynomial, and
$y_3^2=x_3^3-Nx_3/(48D)+M/(864D^2)$.
The Jacobian $J_Y(\mathbb Q)$ has rank zero.
\end{lemma}
\begin{proof}
The two displayed transformations commute and have order $3$.
The first cyclically permutes the three nonzero points of order $2$.
The second fixes the full $2$-structure and multiplies the point
of order $7$ by $4$. They generate a group of order $9$.

We give the descent calculation. The Tate normal form for a point
of order $7$ has
\[
 b_7=v^2(v-1),\qquad c_7=v(v-1).
\]
Put $H=v^2-v+1$, $s=(v^3-3v+1)/(v(v-1))$, and $t=s-8$.
The diamond operator of order $3$ sends $v$ to $(v-1)/v$.
Completing the square and depressing the cubic, followed by
scaling by $H$, gives $E_u$ after adjoining $u^2=t$.
The exact invariant identities are
\[
 D=\frac{H^3}{v^2(v-1)^2},\quad
 \frac{c_4}{H^4}=\frac ND,\quad
 \frac{c_6}{H^6}=-\frac M{D^2},\quad
 \frac{\Delta}{H^{12}}=\frac{t}{D^4}.
\]
The square root $u$ records the orientation of the full
$2$-structure, and is fixed by its alternating subgroup.
The point $14Q$ supplies the nonzero point of order $3$.
These functions are invariant under $G_9$.

The resulting field has degree $8$ over $\mathbb Q(u)$.
The two successive equations have degrees $4$ and $2$ and remain
irreducible over $\mathbb F_5(u)$.
Moreover
\[
 j=\frac{N^3D}{t}
\]
has degree $16$ as a function of $u$. Hence the field has degree
$128$ over $\mathbb Q(j)$. The projective index for $X$ is $1152$,
so this is the full invariant field of $G_9$.
The quotient has good reduction at $5$, since $X$ has good
reduction there and $\#G_9=9$ is prime to $5$.
Its normalized finite function field has genus $5$.
Finally $J_Y$ is a quotient of $J_X$, whose rational rank is zero
by \cite[Theorem~3.1]{DEHMZ}.
\end{proof}

\begin{lemma}\label{lem:decic242-quotient5-hecke}
Normalize $J_Y$ at either rational cusp. The reduction image
of $J_Y(\mathbb Q)$ at $5$ is contained in a subgroup
\[
 K\simeq C_2\times C_2\times C_2\times C_{24}
 \subset J_Y(\mathbb F_5),
\]
of index $18$.
\end{lemma}
\begin{proof}
The curve $Y$ has eight geometric cusps. Their fields are
\[
 \mathbb Q,\quad\mathbb Q,\quad\mathbb Q(\sqrt{21}),\quad
 \mathbb Q(\sqrt{-7}),\quad\mathbb Q(\sqrt{-3}),
\]
where each quadratic field contributes two geometric cusps.
To distinguish the branches, put
\[
 \rho=\frac{M}{12ND},\qquad h=\frac{y_3}{x_3-\rho}.
\]
At $u=0$, the two branch pairs have residues
$h^2=1/84$ and $h^2=-1/28$.
At $u=\infty$, they have residues
$h^2=1/4$ and $h^2=-1/12$.
The first three quadratic fields and the two rational cusps follow.
At $5$ the cusp signature is $1^4 2^2$.

The proven finite Jacobian is
\[
 J_Y(\mathbb F_5)=[2,6,6,48].
\]
The finite cusps generate its entire $2$-primary subgroup,
with invariants $[2,2,2,16]$.
The global rational cuspidal divisors have image $[2,2,8]$.
We do not assume that this is the entire rational image.

On $Y$, the diamond operator $\langle11\rangle$ sends
$(u,x_3,y_3)$ to $(u,x_3,-y_3)$.
The operator
\[
 A_{11}=T_{11}-1-11\langle11\rangle
\]
annihilates $J_Y(\mathbb Q)$.
Indeed the Jacobian has rank zero. Its rational torsion injects
at $5$, and has no $11$-primary part since
$\#J_Y(\mathbb F_5)=3456$. Reduction at $11$ is therefore
injective on this torsion, and the Eichler--Shimura relation gives
the assertion.

We compute $A_{11}$ on the finite cusps using their exact global
Frobenius and diamond permutations. Three ordinary closed cubic
places, together with the finite cusps, generate
$J_Y(\mathbb F_5)$. For each cubic place we lift the level
structure over $\mathbb F_{5^9}$ and compute the twelve cyclic
$11$-isogenies. Equal projected closed places are combined before
dividing the norm divisor by $3$; every resulting multiplicity
is checked to be integral.
The images respect every relation of the full finite group.
The exact kernel is $[2,2,2,24]$, and the image is $[3,6]$.
Independent Sage calculations reconstruct all $36$ checked
kernel-polynomial isogenies and verify that their norms exhaust
the respective $11$-division polynomials. They also reconstruct
the full operator and its kernel by integer linear algebra.
Thus the rational image is contained in $K=\ker A_{11}$.
\end{proof}

